\begin{proof}[Proof of \cref{obj:thm:welfare-identity}]
1.  From the well-formedness of the observed law, extract the measurability of
\(\tau_P\) and the identity \(\tau_P(x)=\mu_1(x)-\mu_0(x)\).  By
\cref{obj:ass:bounded-outcome}, \(|\mu_0(x)|\le1\) and
\(|\mu_1(x)|\le1\) for every \(x\).  Hence, for every \(x\),
\[
|\tau_P(x)|
=
|\mu_1(x)-\mu_0(x)|
\le |\mu_1(x)|+|\mu_0(x)|
\le 2 .
\]
% lean: regret_eq_disagreement_integral

2.  The set \(\{x:0\le\tau_P(x)\}\) is measurable because \(\tau_P\) is
measurable, and the set \(\{x:\pi(x)=1\}\) is measurable because \(\pi\) is
measurable.  Therefore the two integrands
\[
x\mapsto
\begin{cases}
\tau_P(x),&0\le \tau_P(x),\\
0,&\text{otherwise},
\end{cases}
\qquad
x\mapsto
\begin{cases}
\tau_P(x),&\pi(x)=1,\\
0,&\text{otherwise},
\end{cases}
\]
are measurable and, by the preceding bound \(|\tau_P|\le2\), integrable with
respect to the probability measure \(P_X\).  This is the integrability bookkeeping needed to
combine the two welfare integrals into one integral of a difference.
% lean: regret_eq_disagreement_integral

3.  By \cref{obj:def:welfare-regret}, \(V_P(\pi)=\int\mathbf 1\{\pi(x)=1\}\tau_P(x)\,dP_X(x)\) and \(\pi^\star_P(x)=\mathbf 1\{\tau_P(x)\ge0\}\), so \(\mathbf 1\{\pi^\star_P(x)=1\}=\mathbf 1\{0\le\tau_P(x)\}\).  Expanding \(R_P(\pi)=V_P(\pi^\star_P)-V_P(\pi)\) therefore gives
\[
R_P(\pi)
=
\int
\left[
\mathbf 1\{0\le\tau_P(x)\}\tau_P(x)
-
\mathbf 1\{\pi(x)=1\}\tau_P(x)
\right]\,dP_X(x).
\]
The integrability of Step 2 justifies writing the difference of the two integrals as the integral of the pointwise
difference.
% lean: regret_eq_disagreement_integral

4.  It remains only to identify the pointwise difference.  Fix \(x\).  If
\(\tau_P(x)\ge0\), then \(\pi^\star_P(x)=1\).  When \(\pi(x)=1\), both terms equal \(\tau_P(x)\) and the difference is
zero, and indeed there is no disagreement; when \(\pi(x)\ne1\), the difference is \(\tau_P(x)=|\tau_P(x)|\), and there is
disagreement.  If \(\tau_P(x)<0\), then \(\pi^\star_P(x)=0\), so the first term is zero.
When \(\pi(x)=1\), the difference is \(-\tau_P(x)=|\tau_P(x)|\), and there is disagreement; when
\(\pi(x)\ne1\), both terms are zero and there is no disagreement.  Thus, for every \(x\),
\[
\mathbf 1\{0\le\tau_P(x)\}\tau_P(x)
-
\mathbf 1\{\pi(x)=1\}\tau_P(x)
=
|\tau_P(x)|\,\mathbf 1\{\pi(x)\ne\pi^\star_P(x)\}.
\]
% lean: regret_eq_disagreement_integral

5.  Substituting this pointwise identity into the integral yields
\[
R_P(\pi)
=
\int_{\mathcal X}
|\tau_P(x)|\,\mathbf 1\{\pi(x)\ne\pi^\star_P(x)\}\,dP_X(x)
=
E_P\!\left[|\tau_P(X)|\,\mathbf 1\{\pi(X)\ne\pi^\star_P(X)\}\right],
\]
the second equality because \(P_X\) is the covariate marginal of \(P\) and the integrand is a function of \(X\) alone.  This is the disagreement-integral representation of the regret.
% lean: regret_eq_disagreement_integral
\end{proof}

\begin{proof}[Proof of \cref{obj:thm:margin-localization}]
\begin{enumerate}
\item \emph{The constant.}  Fix
\[
C:=|C_m|+2+u_0^{-\alpha}.
\]
\Cref{obj:ass:margin} gives \(0\le\alpha\), \(C_m>0\), and \(u_0>0\), so every summand is nonnegative and \(C>0\); moreover
\[
C\ge C_m+1
\qquad\text{and}\qquad
C\ge u_0^{-\alpha}.
\]
This constant depends only on \(C_m\), \(u_0\), and \(\alpha\).  It remains to prove the displayed inequality for an arbitrary law \(P\) and every \(\pi\in\Pi\). % lean: margin_localization

\item \emph{Two elementary facts.}  Let \(r=R_P(\pi)\) and \(D_\pi=\{x:\pi(x)\ne\pi^\star_P(x)\}\) as in \cref{obj:def:disagreement}.  By \cref{obj:thm:welfare-identity},
\[
r=\int_{\mathcal X} |\tau_P(x)|\,\mathbf 1_{D_\pi}(x)\,dP_X(x),
\]
whose integrand is nonnegative, so \(r\ge0\); the integral is finite because \cref{obj:ass:bounded-outcome} gives \(|\tau_P|=|\mu_1-\mu_0|\le2\).  Also \(P_X(D_\pi)\le1\), since \(P_X\) is a probability measure. % lean: regret_eq_disagreement_integral

\item \emph{The basic decomposition.}  For every \(u\in(0,u_0]\),
\[
P_X(D_\pi)\le C_m u^\alpha+\frac{r}{u}.
\]
Indeed, write
\[
Z=\{x:\tau_P(x)=0,\ x\in D_\pi\},\qquad
S_u=\{x:0<|\tau_P(x)|\le u\},\qquad
B_u=D_\pi\cap\{x:u<|\tau_P(x)|\}.
\]
Every \(x\in D_\pi\) lies in one of these three sets, according to whether \(\tau_P(x)=0\), \(0<|\tau_P(x)|\le u\), or \(|\tau_P(x)|>u\); hence
\[
D_\pi\subseteq Z\cup S_u\cup B_u .
\]
\Cref{obj:ass:zero-effect} gives \(P_X(Z)=0\): in its first alternative \(P_X\{\tau_P=0\}=0\), so its subset \(Z\) is null; in its second alternative, applied to \(\pi\in\Pi\), the set \(Z\) of zero-contrast disagreements is null by hypothesis.  \Cref{obj:ass:margin}, applied at the admissible window \(u\le u_0\), gives
\[
P_X(S_u)\le C_m u^\alpha .
\]
Finally, on \(B_u\) the integrand of the regret identity is at least \(u\) (since \(B_u\subseteq D_\pi\) and \(|\tau_P|>u\) there), and the integrand is nonnegative, so restricting the integral to \(B_u\) gives
\[
u\,P_X(B_u)\le \int_{B_u}|\tau_P(x)|\,\mathbf 1_{D_\pi}(x)\,dP_X(x)\le r,
\qquad\text{hence}\qquad
P_X(B_u)\le \frac ru .
\]
Monotonicity and subadditivity of \(P_X\) along the displayed inclusion now give the asserted decomposition.  The measurability of \(Z\), \(S_u\), \(B_u\) and the finiteness of the integrals are the standard bookkeeping supplied by the well-formedness and bounded-outcome hypotheses. % lean: disagreement_measure_le_margin_plus_regret_over_u

\item \emph{The case \(\alpha=0\).}  Then \(\alpha/(1+\alpha)=0\) and \(r^0=1\), while \(u_0^{-0}=1\) gives \(C=C_m+3\ge1\).  Hence
\[
P_X(D_\pi)\le1\le C=C\,r^{\alpha/(1+\alpha)} .
\]
% lean: margin_localization

\item \emph{The case \(\alpha>0\), \(r>0\), small window.}  Set
\[
u=r^{1/(1+\alpha)}>0,
\]
and suppose first \(u\le u_0\).  Then the two exponent identities
\[
u^\alpha=r^{\alpha/(1+\alpha)},
\qquad
\frac ru=r^{1-1/(1+\alpha)}=r^{\alpha/(1+\alpha)}
\]
hold, so Step 3 gives
\[
P_X(D_\pi)
\le C_m u^\alpha+\frac r u
=(C_m+1)\,r^{\alpha/(1+\alpha)}
\le C\,r^{\alpha/(1+\alpha)},
\]
the last step by \(C_m+1\le C\) and \(r^{\alpha/(1+\alpha)}\ge0\). % lean: margin_localization

\item \emph{The case \(\alpha>0\), \(r>0\), large window.}  Suppose instead \(u>u_0\).  Since \(\alpha>0\), raising to the power \(\alpha\) is increasing, so
\[
u_0^\alpha\le u^\alpha=r^{\alpha/(1+\alpha)}.
\]
Multiplying by \(u_0^{-\alpha}>0\) and using \(u_0^{-\alpha}\le C\),
\[
1=u_0^{-\alpha}u_0^\alpha
\le u_0^{-\alpha}\,r^{\alpha/(1+\alpha)}
\le C\,r^{\alpha/(1+\alpha)} .
\]
Combining with \(P_X(D_\pi)\le1\) from Step 2 gives the desired bound in this subcase. % lean: margin_localization

\item \emph{The case \(\alpha>0\), \(r=0\).}  Let \(\varepsilon>0\), put
\[
a=\frac{\varepsilon}{C_m+1}>0,
\qquad
u=\min\{u_0,a^{1/\alpha}\}.
\]
Then \(0<u\le u_0\) and \(u^\alpha\le a\), so Step 3 with \(r=0\) gives
\[
P_X(D_\pi)\le C_m u^\alpha+\frac 0u
=C_m u^\alpha
\le C_m a
=\frac{C_m}{C_m+1}\,\varepsilon
\le \varepsilon.
\]
Since \(\varepsilon>0\) was arbitrary, \(P_X(D_\pi)=0\).  Also \(\alpha/(1+\alpha)>0\), so
\[
r^{\alpha/(1+\alpha)}=0^{\alpha/(1+\alpha)}=0,
\]
and therefore
\[
P_X(D_\pi)=0\le C\,r^{\alpha/(1+\alpha)} .
\]
% lean: margin_localization
\end{enumerate}
\end{proof}

\begin{proof}[Proof of \cref{obj:thm:minimax-lower}]
1. \emph{The three ingredients and the constant.}  By \cref{obj:lem:regret-separation}, there are constants \(c_{\mathrm{sep}}>0\) and \(N_1\) such that, for all \(n\ge N_1\) and every \(\pi\in\Pi\),
\[
\max_{\sigma\in\{+,-\}} R_{P_{n,\sigma}}(\pi)
\ge
c_{\mathrm{sep}}\,h_n^{1+\alpha}.
\]
By \cref{obj:lem:two-point-divergence}, there are \(C_\chi>0\) and \(N_2\) such that, for all \(n\ge N_2\), the product laws satisfy \(P_{n,+}^{\otimes n}\ll P_{n,-}^{\otimes n}\) with square-integrable density deviation and
\[
\chi^2(P_{n,+}^{\otimes n}\,\|\,P_{n,-}^{\otimes n})\le C_\chi .
\]
Applying \cref{obj:lem:le-cam-two-point-chisq} to this fixed budget gives a constant \(c_{\mathrm{test}}=c_{\mathrm{test}}(C_\chi)>0\) such that every measurable event \(A\) in the \(n\)-sample space satisfies
\[
P_{n,+}^{\otimes n}(A^c)+P_{n,-}^{\otimes n}(A)\ge c_{\mathrm{test}} .
\]
Set
\[
c=\frac{c_{\mathrm{sep}}\,c_{\mathrm{test}}}{2}>0 .
\]
% lean: minimax_lower

2. \emph{Membership and the rate identity.}  By \cref{obj:lem:witness-membership}, enlarging \(N_1,N_2\) if necessary, both witness laws \(P_{n,+}\) and \(P_{n,-}\) belong to \(\mathcal P_{\alpha,\gamma}\) for all \(n\) under consideration.  Since
\[
h_n=n^{-1/(2+\alpha+\beta_{\alpha,\gamma})},
\qquad
r_\star(\alpha,\gamma)=\frac{1+\alpha}{2+\alpha+\beta_{\alpha,\gamma}},
\]
raising \(h_n\) to the power \(1+\alpha\) gives the exact identity
\[
h_n^{1+\alpha}=n^{-r_\star(\alpha,\gamma)} .
\]
The policy class is nonempty, so the collection of measurable \(\Pi\)-valued estimators is nonempty (a constant estimator qualifies) and the infimum in \cref{obj:def:minimax-regret} is over a nonempty set; and regret is nonnegative and bounded by \(2\) under \cref{obj:ass:bounded-outcome}, so all displayed expectations are finite.  This is the only role of the nonemptiness and boundedness bookkeeping. % lean: witness_membership % lean: lawRegret_nonneg_of_wellformed % lean: lawRegret_le_two_of_wellformed

3. \emph{The threshold event.}  Fix such an \(n\), and let \(\widehat\pi\) be any measurable \(\Pi\)-valued estimator.  Write
\[
R_+
=
\mathbb E_{P_{n,+}^{\otimes n}} R_{P_{n,+}}(\widehat\pi),
\qquad
R_-
=
\mathbb E_{P_{n,-}^{\otimes n}} R_{P_{n,-}}(\widehat\pi),
\]
put
\[
s_n=c_{\mathrm{sep}}h_n^{1+\alpha}>0,
\qquad
A=\left\{\text{samples}:\ s_n\le R_{P_{n,-}}(\widehat\pi)\right\}.
\]
The set \(A\) is measurable because the map from the sample to \(R_{P_{n,-}}(\widehat\pi)\) is measurable, which is part of the estimator condition in \cref{obj:def:minimax-regret}. % lean: minimax_lower

4. \emph{Two Markov-type threshold bounds.}  For a probability measure \(Q\), a nonnegative \(Q\)-integrable function \(f\), and a level \(s>0\), the pointwise inequality \(f\ge s\,\mathbf 1\{f\ge s\}\) integrates to
\[
s\,Q\{f\ge s\}\le\int f\,dQ .
\]
Applying this to the nonnegative integrable regret maps under the two product laws,
\[
s_n\,P_{n,-}^{\otimes n}(A)\le R_-,
\qquad
s_n\,P_{n,+}^{\otimes n}\left\{s_n\le R_{P_{n,+}}(\widehat\pi)\right\}\le R_+ .
\]
% lean: mul_meas_ge_le_integral_of_nonneg
On \(A^c\) we have \(R_{P_{n,-}}(\widehat\pi)<s_n\); applying the separation bound of Step 1 pointwise to the realized policy \(\widehat\pi\in\Pi\), the maximum of the two regrets is at least \(s_n\), so it must be the plus-regret that attains it.  Hence
\[
A^c\subseteq \left\{s_n\le R_{P_{n,+}}(\widehat\pi)\right\},
\]
and by monotonicity of \(P_{n,+}^{\otimes n}\) together with the second threshold bound,
\[
s_n\,P_{n,+}^{\otimes n}(A^c)\le R_+ .
\]

5. \emph{Combining with the testing floor.}  Adding the two displayed bounds and inserting the testing lower bound of Step 1,
\[
s_n\,c_{\mathrm{test}}
\le
s_n\left\{P_{n,+}^{\otimes n}(A^c)+P_{n,-}^{\otimes n}(A)\right\}
\le
R_+ + R_- .
\]
Since \(R_++R_-\le 2\max\{R_+,R_-\}\), dividing by \(2\) gives
\[
\max\{R_+,R_-\}
\ge
\frac{s_n c_{\mathrm{test}}}{2}
=
\frac{c_{\mathrm{sep}}c_{\mathrm{test}}}{2}\,h_n^{1+\alpha}
=
c\,h_n^{1+\alpha}
=
c\,n^{-r_\star(\alpha,\gamma)} .
\]
% lean: le_cam_two_point_chisq

6. \emph{Passing to the minimax risk.}  Because \(P_{n,+},P_{n,-}\in\mathcal P_{\alpha,\gamma}\) by Step 2, each of \(R_+\) and \(R_-\) is one of the values over which the supremum in \cref{obj:def:minimax-regret} is taken, so
\[
\max\{R_+,R_-\}
\le
\sup_{P\in\mathcal P_{\alpha,\gamma}}
\mathbb E_{P^{\otimes n}} R_P(\widehat\pi).
\]
Therefore every admissible estimator satisfies
\[
\sup_{P\in\mathcal P_{\alpha,\gamma}}
\mathbb E_{P^{\otimes n}} R_P(\widehat\pi)
\ge
c\,n^{-r_\star(\alpha,\gamma)},
\]
with \(c\) not depending on \(\widehat\pi\) or on \(n\).  Taking the infimum over all measurable \(\Pi\)-valued estimators in \cref{obj:def:minimax-regret} gives
\[
M_n(\alpha,\gamma)\ge c\,n^{-r_\star(\alpha,\gamma)}
\]
for all sufficiently large \(n\). % lean: minimax_lower
\end{proof}

\begin{proof}[Proof of \cref{obj:oeq:feasible-upper}]
\begin{enumerate}
\item \emph{The schedules.}  The clipping and localization schedules are those of \cref{obj:def:feasible-rate}: for \(\gamma>0\),
\[
q_n=q_0n^{-s_{\mathrm{feas}}},\qquad
u_n=\bar u\,n^{-t_{\mathrm{feas}}},
\]
while for \(\gamma=0\) the clipping schedule is fixed, \(q_n=q_0\) for every \(n\).  The input restrictions assumed in the theorem give \(q_0>0\) in both regimes; when \(\gamma>0\) they also give \(\bar u>0\), \(q_0\le\min\{1/2,c_o\bar u^\gamma\}\), and — this is the admissibility recorded in \cref{obj:def:feasible-rate} — the eventual inequality
\[
q_n\le c_o u_n^\gamma .
\]
In particular \(q_n>0\) for all sufficiently large \(n\). % lean: qSched % lean: feasibleAdmissible

\item \emph{The exceptional empty-domain case.}  Suppose \(\gamma=0\) and no law satisfies the law-class conditions of \cref{obj:def:law-class}.  Then the set over which the supremum in \cref{obj:def:upper-risk} is taken is empty, so \(U_n(\alpha,\gamma,a,c;\widehat\eta)=0\), while the right-hand side of the asserted bound is nonnegative.  Taking \(C=1\) and \(p=0\) settles this case, and it is excluded from now on. % lean: feasible_upper

\item \emph{The clip lies in the admissible clipping range.}  If \(\gamma=0\), the excluded case supplies some law \(P\in\mathcal P_{\alpha,\gamma}\); its strict-overlap endpoint condition in \cref{obj:def:law-class} gives \(\underline p\le1/2\), and the input restriction \(q_0\le\underline p/2\) then gives
\[
q_n=q_0\le\frac{\underline p}2\le\frac14\le\frac12 .
\]
If \(\gamma>0\), the maximizer of \cref{obj:def:feasible-rate} lies in the feasible box, so \(s_{\mathrm{feas}}\ge0\), hence \(n^{-s_{\mathrm{feas}}}\le1\) for \(n\ge1\), and therefore
\[
q_n=q_0n^{-s_{\mathrm{feas}}}\le q_0\le\frac12 .
\]
% lean: feasibleMaximizer_mem
Likewise, when \(\gamma>0\), \(t_{\mathrm{feas}}\ge0\) gives \(n^{-t_{\mathrm{feas}}}\le1\) for \(n\ge1\), so with \(\bar u\le u_0\),
\[
0<u_n=\bar u\,n^{-t_{\mathrm{feas}}}\le\bar u\le u_0 .
\]

\item \emph{The two ingredients.}  \Cref{obj:lem:crude-localized-master-bound}, applied with these deterministic schedules (whose hypotheses \(0<q_n\le1/2\), and \(q_n\) eventually constant when \(\gamma=0\), were just verified), supplies constants \(C_{\mathrm M}>0\) and \(p_{\mathrm M}\ge0\) such that, for all sufficiently large \(n\), every law \(P\) in the side-condition domain of \cref{obj:def:upper-risk} satisfies the corresponding master inequality.  \Cref{obj:lem:clip-balance-exponent}, applied to the same schedules and the nuisance-rate exponents of \cref{obj:ass:polynomial-nuisance-exponents}, supplies constants \(C_{\mathrm B}>0\) and \(p_{\mathrm B}\ge0\) bounding the deterministic bracket of the master bound by
\[
C_{\mathrm B}n^{-r_{\mathrm{feas}}}(\log n)^{p_{\mathrm B}},
\qquad
r_{\mathrm{feas}}=r_{\mathrm{up}}
\]
by \cref{obj:def:feasible-rate}.  Both constant pairs are chosen before \(n\).  Set
\[
C=C_{\mathrm M}C_{\mathrm B}>0,
\qquad
p=p_{\mathrm M}+p_{\mathrm B}\ge0 .
\]
% lean: crude_localized_master_bound % lean: clip_balance_exponent

\item \emph{The case \(\gamma=0\).}  For every \(P\) in the domain of \cref{obj:def:upper-risk} and all sufficiently large \(n\), the strict-overlap branch of \cref{obj:lem:crude-localized-master-bound} applies because \(q_n=q_0\le\underline p/2\), and gives
\[
\mathbb E_P R_P(\widehat\pi_n)
\le
C_{\mathrm M}\left\{n^{-A_\alpha}+r_{\mu,n}r_{e,n}\right\}(\log n)^{p_{\mathrm M}} .
\]
The \(\gamma=0\) branch of \cref{obj:lem:clip-balance-exponent} bounds the brace by \(C_{\mathrm B}n^{-r_{\mathrm{feas}}}(\log n)^{p_{\mathrm B}}\).  Since \(\log n>0\) for \(n\ge2\), the two logarithmic factors multiply as \((\log n)^{p_{\mathrm B}}(\log n)^{p_{\mathrm M}}=(\log n)^{p}\), and therefore
\[
\mathbb E_P R_P(\widehat\pi_n)
\le
C\,n^{-r_{\mathrm{feas}}}(\log n)^{p} .
\]

\item \emph{The case \(\gamma>0\).}  For all sufficiently large \(n\) we have \(u_n>0\), \(u_n\le u_0\) (item 3), and \(q_n\le c_ou_n^\gamma\) (item 1), so the positive-\(\gamma\) branch of \cref{obj:lem:crude-localized-master-bound} applies and yields, for each \(P\) in the domain,
\[
\mathbb E_P R_P(\widehat\pi_n)
\le
C_{\mathrm M}
\left\{
n^{-r_\star(\alpha,\gamma)}
+(nq_n^2)^{-A_\alpha}
+\frac{r_{\mu,n}r_{e,n}}{q_n}
+r_{\mu,n}u_n^{\alpha/2}q_n^{1/(2\gamma)}
+\frac{r_{\mu,n}^2}{u_n}
\right\}
(\log n)^{p_{\mathrm M}} .
\]
The positive-\(\gamma\) branch of \cref{obj:lem:clip-balance-exponent} bounds the brace by \(C_{\mathrm B}n^{-r_{\mathrm{feas}}}(\log n)^{p_{\mathrm B}}\), and combining the two displays as in the previous item gives
\[
\mathbb E_P R_P(\widehat\pi_n)
\le
C\,n^{-r_{\mathrm{feas}}}(\log n)^{p} .
\]

\item \emph{Passing to the supremum.}  In both cases the bound \(C\,n^{-r_{\mathrm{feas}}}(\log n)^{p}\) does not depend on the particular law \(P\) in the domain of \cref{obj:def:upper-risk}, and it is nonnegative; hence it also bounds the supremum defining the conditional upper risk,
\[
U_n(\alpha,\gamma,a,c;\widehat\eta)
\le
C\,n^{-r_{\mathrm{feas}}}(\log n)^{p}
=
C\,n^{-r_{\mathrm{up}}}(\log n)^{p},
\]
for all sufficiently large \(n\), with the single pair \((C,p)\) fixed in item 4.  This is the claimed bound. % lean: feasible_upper
\end{enumerate}
\end{proof}
